\documentclass[10pt,a4paper]{amsart}
\usepackage[margin=19mm]{geometry}
\usepackage{amsmath,amssymb,mathtools}
\usepackage[scr=rsfs,cal=euler]{mathalfa}
\usepackage{xfrac}
\usepackage[unicode,hidelinks,bookmarks=false]{hyperref}
\newcommand{\ie}{i.e.\ }
\newcommand{\cf}{cf.\ }
\newcommand{\resp}{respectively\ }
\newcommand{\Subsection}{\subsection}
\providecommand{\nobreakdash}{\nobreak\hbox{-}}
\setlength{\emergencystretch}{2em}
\multlinegap0pt
\usepackage[shortlabels]{enumitem}
\usepackage{xcolor}
\usepackage{mathtools}
\newtheorem{lemma}{Lemma}
\newtheorem{proposition}{Proposition}
\newcommand{\E}{\mathbb E}
\newcommand{\R}{\mathbb R}
\newcommand{\abs}[1]{\lvert #1 \rvert}
\newcommand{\eps}{\varepsilon}
\newcommand{\ip}[2]{\langle #1,#2\rangle}
\newcommand{\norm}[1]{\lVert #1 \rVert}
\title{Stable Phase Retrieval for Spans of Independent Random Variables}
\author{Pedro Abdalla, Jaume de Dios Pont, Jo\~ao P. G. Ramos, Mitchell A. Taylor}
\date{Source: arXiv:2607.06693v1 (CC BY 4.0)}
\begin{document}
\maketitle

\noindent\emph{Source and licence.} Stable Phase Retrieval for Spans of Independent Random Variables. Version arXiv:2607.06693v1 (CC BY 4.0), distributed by arXiv under the Creative Commons Attribution 4.0 International licence, http://creativecommons.org/licenses/by/4.0/, as recorded in the arXiv licence metadata for this exact version and shown on the abstract page https://arxiv.org/abs/2607.06693v1. Full licence notice: \texttt{licenses/arxiv-2607.06693-CC-BY-4.0.txt}.\par\smallskip

\noindent\textit{Editorial scope.} Counted unit: the article's proposition prop:concentrated (source0.tex lines 1472--1481). The article's complete original proof of it is delivered with it (source0.tex lines 1483--1822, one \texttt{proof} environment of 340 lines). No mathematics is written, repaired or re-derived: the statement and the proof are the article's own lines, in the article's own notation, and no retained line is substituted or re-flowed.  Retained uncounted as the source-local context the proof needs: lines 397--411; lines 949--961; lines 1402--1419; lines 1458--1463.  Nothing was omitted from any retained span, so the package carries no omission record.  No retained line contains a citation, so the wrapper carries no bibliography and no bibliography entry.  No retained line contains a figure environment, an image-inclusion command or drawing code, so no image is delivered and the package declares no asset dependency.  Editorial formatting only, in addition: an AMS \texttt{amsart} preamble that restates the article's own declarations and macros used by the retained lines, namely the environments \texttt{lemma}, \texttt{proposition} on separate counters, together with the standard packages those lines need.  The retained environments print in this document as Lemma 1, Proposition 1 in order of appearance, whereas the article prints the same items with the numbering of its own full document.  The complete native source of the article as distributed in the arXiv e-print for this exact version is bundled unchanged as \texttt{sources/204655/source0.tex}.
\begin{lemma}[$L^1$ lower bound for independent sums]\label{lem:l1lower}
Let $\zeta_1,\dots,\zeta_N$ be independent, centered, real-valued random variables and
assume that
\[
\norm{\zeta_i}_{L^1}\ge \delta \norm{\zeta_i}_{L^2}>0
\qquad(1\le i\le N).
\]
Then for every scalar family $(c_i)_{i=1}^N$, we have
\[
\norm{\sum_{i=1}^N c_i\zeta_i}_{L^1}
\ge
\frac{\delta}{2\sqrt2}
\Bigl(\sum_{i=1}^N c_i^2\norm{\zeta_i}_{L^2}^2\Bigr)^{1/2}.
\]
\end{lemma}

\begin{lemma}[Identification of the head profile]\label{lem:head-sign}
Let the notation be as above. Then either there exists $\tau\in\{\pm1\}$ such that
\[
a=\tau b
\qquad\text{and}\qquad
\norm{a_{\mathrm{head}}^{(n)}-\tau b_{\mathrm{head}}^{(n)}}_2\to0,
\]
or else $\eta$ is a Rademacher variable and
\[
a_i=b_i=0
\qquad(i\ge2).
\]
\end{lemma}

\begin{proposition}
\label{prop:head-entry}
The following estimates hold.
\begin{enumerate}[label=\rm(\arabic*)]
\item For any indices $i,j\in H\cap\{2,3,\dots\}$,
\[
|a_i a_j-b_i b_j|\le K_{\delta,\eta}\,\eps.
\]
\item For any index $j\in H\cap\{2,3,\dots\}$,
\[
|a_1a_j-b_1b_j|\le K_{\delta,\eta}\,\eps.
\]
\item If $\eta^2\not\equiv1$, then
\[
|a_1^2-b_1^2|\le K_{\delta,\eta}\,\eps.
\]
\end{enumerate}
\end{proposition}

\begin{lemma}\label{lem:rank-one}
For $x,y\in\R^m$, let $A(x,y)_{ij}=x_ix_j-y_iy_j$ be the difference matrix between $x$ and $y$. Then
\[
\norm{x-y}_2^2\,\norm{x+y}_2^2\le 2\norm{A(x,y)}_F^2.
\]
\end{lemma}

\begin{proposition}[Concentrated case]\label{prop:concentrated}
If there exists $\tau\in\{\pm 1\}$ such that
\[
\norm{(a-\tau b)\mathbf 1_{H^c}}_2\le \frac{\theta}{32},
\]
then there exists a constant $c_{\mathrm{line}}(\delta,\eta)>0$ such that
\[
\eps\ge c_{\mathrm{line}}(\delta,\eta).
\]
\end{proposition}

\begin{proof}
We distinguish two cases.
\vspace{0.5em}

\noindent\emph{Case 1: $\eta$ is not Rademacher or $H\cap\{2,3,\dots\}$ is
nonempty.}
In this case, the head coefficients are linked by
Proposition~\ref{prop:head-entry}. %either directly through the diagonal
%estimate for the exceptional coordinate or, in the Rademacher case, through a
%good head coordinate exactly as in the proof of Lemma~\ref{lem:head-sign}.  
In particular, by direct calculation and an argument similar to Lemma~\ref{lem:head-sign} in the Rademacher case, 
there exists a sign, which must be the same $\tau$, such that the
difference matrix of $a\mathbf 1_H$ and $\tau b\mathbf 1_H$ has all
entries bounded by $K_{\delta,\eta}\eps$. Thus, applying Lemma~\ref{lem:rank-one} with
\[
x:=a\mathbf 1_H,
\qquad
y:=\tau b\mathbf 1_H,
\]
we obtain 
\[
\norm{x-y}_2^2\,\norm{x+y}_2^2
\le 2\#H^2 K_{\delta,\eta}^2\eps^2.
\]
Since $(a-\tau b)\mathbf{1}_{H^c}$ has norm at most $\theta/32$ and
$\norm{a-\tau b}_2^2=2$ by orthogonality between $a$ and $b$, we have
\[
\norm{x-y}_2^2
=2-\norm{(a-\tau b)\mathbf 1_{H^c}}_2^2
\ge 1.
\]
	Hence,
	\[
	\norm{x+y}_2\le \sqrt{2}(\#H) K_{\delta,\eta}\eps
	\le \sqrt{2}d^*(\delta,\eta)K_{\delta,\eta}\eps.
	\]
	Set
	\[
	\alpha:=\norm{(a+\tau b)\mathbf 1_H}_2,
	\qquad
	\beta:=\norm{(a-\tau b)\mathbf 1_{H^c}}_2,
	\]
	and note that
	\[
	\alpha\le \sqrt2\,d^*(\delta,\eta)K_{\delta,\eta}\eps,
	\qquad
	\beta\le \theta/32.
	\]
	We now describe the gluing step.  Write
	\[
	p:=(a-\tau b)\mathbf 1_H,
	\qquad
	q:=(a-\tau b)\mathbf 1_{H^c},
	\]
	\[
	r:=(a+\tau b)\mathbf 1_H,
	\qquad
	s:=(a+\tau b)\mathbf 1_{H^c},
	\]
	and let
	\[
	P:=p_1\eta+\sum_{i\in H\cap\{2,3,\dots\}} p_i\xi_i,
	\qquad
	Q:=\sum_{i\notin H} q_i\xi_i,
	\]
	\[
	R:=r_1\eta+\sum_{i\in H\cap\{2,3,\dots\}} r_i\xi_i,
	\qquad
	S:=\sum_{i\notin H} s_i\xi_i.
	\]
	We may decompose
	\[
	X-\tau Y=P+Q,
	\qquad
	X+\tau Y=R+S,
	\]
	so that
	\[
	X^2-Y^2
	=(X-\tau Y)(X+\tau Y)
	=(P+Q)(R+S).
	\]
	Since the supports of $p,r$ and $q,s$ are disjoint, the pair $(P,R)$ is
	independent of the pair $(Q,S)$. In particular, $P$ and $S$ are independent.
	In view of the inequality
	\[
	\abs{PS}\le \abs{X^2-Y^2}+\abs{PR}+\abs{QS}+\abs{QR},
	\]
	 we deduce that
	\begin{equation}
    \label{eq:aux_conc_branch}
	\E[\abs{P}]\,\E[\abs{S}]
	=\E[\abs{PS}]
	\le
	\eps+\E[\abs{PR}]+\E[\abs{QS}]+\E[\abs{QR}].
	\end{equation}
	By Cauchy--Schwarz, we also have
	\[
	\E[\abs{PR}]\le \norm{P}_{L^2}\norm{R}_{L^2},
	\qquad
	\E[\abs{QS}]\le \norm{Q}_{L^2}\norm{S}_{L^2},
	\qquad
	\E[\abs{QR}]\le \norm{Q}_{L^2}\norm{R}_{L^2}.
	\]
	Moreover, we have the identities
	\[
	\norm{P}_{L^2}=\norm{p}_2=\sqrt{2-\beta^2},
	\qquad
	\norm{R}_{L^2}=\norm{r}_2=\alpha,
	\]
	\[
	\norm{Q}_{L^2}=\norm{q}_2=\beta,
	\qquad
	\norm{S}_{L^2}=\norm{s}_2=\sqrt{2-\alpha^2}.
	\]
	Now assume, for the sake of contradiction, that $\eps<c_{\mathrm{line}}(\delta,\eta)$, where
	$c_{\mathrm{line}}(\delta,\eta)$ is chosen so small that
	\[
	\sqrt2\,d^*(\delta,\eta)K_{\delta,\eta}c_{\mathrm{line}}(\delta,\eta)\le1.
	\]
	Under these assumptions, we may conclude that $\alpha\le1$ and
	\[
    1\le \min\{\|P\|_{L^2},\|S\|_{L^2}\}\le \max\{\|P\|_{L^2},\|S\|_{L^2}\} 
	\le \sqrt{2}.
	\]
	Consequently, from \eqref{eq:aux_conc_branch} we may bound
	\begin{equation}
    \label{eq:aux2_conc_branch}
	\E[\abs{P}]\,\E[\abs{S}]
	\le
	\eps+\sqrt2\,\alpha+\sqrt2\,\beta+\alpha\beta.
	\end{equation}
	To lower bound the left-hand side of \eqref{eq:aux2_conc_branch}, we decompose $P$ as 
	\[
	P=p_1\eta+G,
	\qquad \text{where} \quad 
	G:=\sum_{i\in H\cap\{2,3,\dots\}} p_i\xi_i.
	\]
	Note that either $\abs{p_1}\ge \norm{p}_2/\sqrt2$ or $\norm{G}_{L^2}\ge \norm{p}_2/\sqrt2$.
	In the first case, we apply Jensen's inequality to estimate
	\[
	\E[\abs{P}]
	=\E\bigl[\abs{p_1\eta+G}\bigr]
	\ge
	\abs{p_1}\,\E[\abs{\eta}]
	\ge
	\kappa_{\delta,\eta}\norm{p}_2.
	\]
	In the second case, we let $G'$ be an independent copy of $G$ and compute that
	\[
	2\E[\abs{P}]
	=\E\bigl[\abs{p_1\eta+G}+\abs{p_1\eta+G'}\bigr]
	\ge \E[\abs{G-G'}].
	\]
	Then, we write
	\[
	G-G'
	=
	\sum_{i\in H\cap\{2,3,\dots\}} p_i(\xi_i-\xi_i'),
	\]
	where $(\xi_i')$ is an independent copy of $(\xi_i)$.  The variables
	$\xi_i-\xi_i'$ are independent, centered, satisfy
	\[
	\norm{\xi_i-\xi_i'}_{L^2}=\sqrt2,
	\qquad
	\norm{\xi_i-\xi_i'}_{L^1}\ge \delta,
	\]
	and therefore also satisfy
	\[
	\norm{\xi_i-\xi_i'}_{L^1}
	\ge \frac{\delta}{\sqrt2}\norm{\xi_i-\xi_i'}_{L^2}.
	\]
	Applying Lemma~\ref{lem:l1lower}, we may bound
	\[
	\E[\abs{G-G'}]
	\ge
	\frac{\delta}{2\sqrt2}
	\Bigl(\sum_{i\in H\cap\{2,3,\dots\}} p_i^2
	\norm{\xi_i-\xi_i'}_{L^2}^2\Bigr)^{1/2}
	=\frac{\delta}{2}\norm{G}_{L^2}
	\ge 2\kappa_{\delta,\eta}\norm{p}_2.
	\]
	In both cases, we obtain
	\[
	\E[\abs{P}]
	\ge \kappa_{\delta,\eta}\norm{p}_2
	\ge \kappa_{\delta,\eta}.
	\]
	Since $S$ is supported on $H^c$, we may apply
	Lemma~\ref{lem:l1lower} to obtain
	\[
	\E[\abs{S}]
	\ge \frac{\delta}{2\sqrt2}\norm{s}_2
	\ge \frac{\delta}{2\sqrt2}.
	\]
	Combining the last two displays with the inequality in \eqref{eq:aux2_conc_branch}, we see that
	\[
	\frac{\delta\kappa_{\delta,\eta}}{2\sqrt2}
	\le
	\eps+\sqrt2\,\alpha+\sqrt2\,\beta+\alpha\beta.
	\]
	Using the bounds for $\alpha$ and $\beta$, we infer that
	\[
	\frac{\delta\kappa_{\delta,\eta}}{2\sqrt2}
	\le
	\Bigl(1+2d^*(\delta,\eta)K_{\delta,\eta}
	+\frac{\theta}{16}d^*(\delta,\eta)K_{\delta,\eta}\Bigr)\eps
	+\frac{\sqrt2\,\theta}{32}.
	\]
	By the choice of $\theta$, we deduce the inequality
	\[
	\frac{\sqrt2\,\theta}{32}
	\le \frac{\delta\kappa_{\delta,\eta}}{8\sqrt2}.
	\]
	After decreasing $c_{\mathrm{line}}(\delta,\eta)$ if necessary so that
	\[
	\Bigl(1+2d^*(\delta,\eta)K_{\delta,\eta}
	+\frac{\theta}{16}d^*(\delta,\eta)K_{\delta,\eta}\Bigr)
	c_{\mathrm{line}}(\delta,\eta)
	\le \frac{\delta\kappa_{\delta,\eta}}{8\sqrt2}
	\]
	(which is strictly smaller than
$\delta\kappa_{\delta,\eta}/(2\sqrt2)$ whenever
	$\eps<c_{\mathrm{line}}(\delta,\eta)$), we reach a contradiction. We conclude, therefore, that
	$\eps\ge c_{\mathrm{line}}(\delta,\eta)$.
\vspace{0.5em}

\noindent\emph{Case 2: $\eta$ is Rademacher and $H=\{1\}$.}
Write
\[
u:=a\mathbf 1_{H^c},
\qquad
v:=b\mathbf 1_{H^c},
\qquad
\rho:=\theta/32,
\]
and let
\[
U:=\sum_{i\ge2} a_i\xi_i,
\qquad
V:=\sum_{i\ge2} b_i\xi_i.
\]
Since $H=\{1\}$, every coordinate entering $U$ and $V$ is good.  The
 assumption of being concentrated states precisely that
\[
\norm{u-\tau v}_2=\norm{U-\tau V}_{L^2}\le \rho.
\]
Set
\[
R:=U-\tau V,
\qquad
S:=U+\tau V,
\qquad
\alpha:=a_1-\tau b_1,
\qquad
\beta:=a_1+\tau b_1
\]
and decompose
\[
X-\tau Y=\alpha\eta+R,
\qquad
X+\tau Y=\beta\eta+S.
\]
Since $\norm{a}_2=\norm{b}_2=1$ and $\ip{a}{b}=0$, we have
\[
\norm{a-\tau b}_2^2=\norm{a+\tau b}_2^2=2.
\]
Therefore,
\[
\alpha^2+\norm{R}_{L^2}^2
=\alpha^2+\norm{u-\tau v}_2^2
=2,
\qquad
\beta^2+\norm{S}_{L^2}^2
=\beta^2+\norm{u+\tau v}_2^2
=2.
\]
In particular,
\[
|\alpha|\ge \sqrt{2-\rho^2}\ge1.
\]
Our next objective is to estimate $\beta$. To this end, we notice that
\[
|\beta|\le |\alpha\beta|
=|a_1^2-b_1^2|
=\bigl|\norm{v}_2^2-\norm{u}_2^2\bigr|
=\bigl|\langle u-\tau v,u+\tau v\rangle\bigr|.
\]
Thus, it follows that
\[
|\beta|
\le \norm{u-\tau v}_2\norm{u+\tau v}_2
\le \sqrt2\,\rho.
\]
Hence, $\norm{S}_{L^2}^2 = \|u+\tau v\|_2^2\ge1$ and Lemma~\ref{lem:l1lower} gives
\[
\norm{S}_{L^1}
=\Bigl\|\sum_{i\ge2}(a_i+\tau b_i)\xi_i\Bigr\|_{L^1}
\ge \frac{\delta}{2\sqrt2}\norm{u+\tau v}_2
\ge \frac{\delta}{2\sqrt2}.
\]
On the other hand, we have
\[
\norm{R}_{L^1}\le \norm{R}_{L^2}\le \rho.
\]
Now, we condition on the $\sigma$-algebra generated by the good coordinates.  Since
$\eta$ is an independent Rademacher variable, we have
\[
X^2-Y^2
=(X-\tau Y)(X+\tau Y)
=(\alpha\eta+R)(\beta\eta+S)
=\alpha\beta+RS+\eta(\alpha S+\beta R).
\]
Therefore,
\begin{align*}
    \E\bigl[\abs{X^2-Y^2}\mid R,S\bigr]
&=\frac12\Bigl(
\abs{\alpha\beta+RS+\alpha S+\beta R}
+
\abs{\alpha\beta+RS-\alpha S-\beta R}
\Bigr)
\\
&=\max\bigl(\abs{\alpha\beta+RS},\abs{\alpha S+\beta R}\bigr)
\ge \abs{\alpha S+\beta R}.
\end{align*}
Taking expectations on both sides, by the law of iterated expectation, we have 
\begin{equation*}
\E[|X^2-Y^2|]\ge \|\alpha S+\beta R\|_{L^1}\ge  |\alpha|\|S\|_{L^1}-|\beta|\|R\|_{L^1}.
\end{equation*}
Recalling the definition of $\varepsilon$ and the estimates for the $L^1$ norms of $S$ and $R$, we conclude that
\[
\eps
\ge |\alpha|\norm{S}_{L^1}-|\beta|\norm{R}_{L^1}
\ge \frac{\delta}{2\sqrt2}-\sqrt2\,\rho^2.
\]
Since $\rho=\theta/32\le \delta^2/256$ and $0<\delta\le1$, the right-hand side
is bounded below by a positive constant depending only on $\delta$ (for
instance by $\delta/4$).  After decreasing
$c_{\mathrm{line}}(\delta,\eta)$ if necessary, this proves Case~2.
\end{proof}

\end{document}
